\section{Renormalization of Effective QCD with ``Heavy Quark'' in Dimensional Regularization}

Let us consider renormalization of the above effective theory with a ``heavy quark''.
We will show that such a theory can be renormalized perturbatively to all orders in perturbation theory,
first in dimensional regularization where power divergences do not exist. The basic argument
is the same as the all-order proof of the renormalization for HQET, first
presented in Ref.~\cite{Bagan:1993zv}.

The standard proof of the renormalization of QCD chooses the covariant gauge $\partial_\mu A^\mu = 0$, implemented
in the functional integrals with a gauge parameter $\xi$. To recover a local theory after gauge
fixing, one has to introduce the color-octet ghost fields $c^a$ and $\bar c^a$. The resulting theory
has a residual global symmetry called BRST symmetry. The Ward-Takahashi identities
from this symmetry are a key ingredient for showing the theory is renormalizable to all-orders
in perturbation theory in the following sense: After absorbing all the infinities into
wave function renormalization factors of quark ($Z_1$), gluon $(Z_3$), and ghost fields $(Z_c)$, and the multiplicative
renormalization of quark masses ($Z_m$), gauge coupling ($Z_g$) and gauge fixing parameter ($Z_\xi$),
Green's functions of the renormalized elementary fields are UV finite.

In HQET where the Lagrangian has a similar form as Eq.~(\ref{Eq:lag}) except that the vector $n^\mu$ is timelike instead of spacelike, it has been shown that to leading order in the heavy quark mass the Green's functions can be renormalized to all-orders in perturbation theory~\cite{Bagan:1993zv}, where a BRST invariance of the HQET Lagrangian was used. In HQET, heavy quarks do not couple to heavy antiquarks, hence no heavy quark loops can occur. This implies that Green's functions without external heavy quark legs can be renormalized in the same way as in QCD. Only those involving external heavy quark legs need to be considered separately, which yield
a new wave function renormalization constant ($Z_Q$). These Green's functions include the heavy-quark two-point function, the heavy-quark-gluon vertex and the heavy-quark-ghost vertex with one insertion of composite operators coming from the BRST transformation of the heavy quark field. All of them can be shown to be UV finite with the help of Ward-Takahashi identities under the given finite number of renormalization factors in the HQET Lagrangian~\cite{Bagan:1993zv}.

Note that the above arguments are valid independent of whether the heavy quark field in the HQET Lagrangian is defined by a timelike or spacelike vector.
Therefore, the statement about renormalizability of the HQET in Ref.~\cite{Bagan:1993zv} shall also be
true for our auxiliary heavy quark Lagrangian with $n^\mu=(0,0,0,1)$ in Eq.~(\ref{Eq:lag}). In the case of light-cone vector, however,
new divergences appear due to light-cone singularities.

Local composite operators in this effective theory can be renormalized using the standard approach. In particular, the heavy-light composite operators, such as
\begin{equation}
             \overline{j}(x) = \overline{\psi}(x)\Gamma Q(x),
\end{equation}
acquires a divergent factor $Z_j$ (Note that $\overline{j}(x)$
now carries a spinor index). In analogy to the Green's functions, one can show that the overall UV divergence of the insertion of $\bar{j}(x)$ into Green's functions is subtracted by $Z_j$ to all orders of perturbation theory. Therefore, we can write $\overline{j}= Z_j \overline{j}_R$ where $\overline{j}_R$
is a renormalized operator.
 The calculation of its anomalous dimension
has been worked out in the HQET~\cite{Politzer:1988wp,Ji:1991pr}.
For the auxiliary heavy quark field, actually one can explicitly verify that at one-loop in dimensional regularization (with $D=4-2\epsilon$)
\begin{equation}\label{eq:hl}
   Z_j = 1+ \frac{\alpha_s}{2\pi\epsilon} \ ,
\end{equation}
which is the same as in HQET.
